\chapter{ Revisão bibliográfica}
\section{Desenvolvimento do método e principais formulações}
%Explicar história do surgimento do (MFDM) 
\indent O método das diferenças finitas miméticas (MFDM) é uma técnica númericas em malhas poligonais que preserva as identidades 
fundamentais do calculo vetorial e as leis da consevação física, diferindo-se dos métodos tradicionais de diferenças finitas que aproximam
as derivadas usando a expansão da série de Taylor em pontos discretos que pode acabar quebrando propriedades matemáticas globais.
No MFDM, as váriaveis são alocadas de acordo com a topologia do problema atribuindo propriedades físicas a entidades geométricas como nós, 
faces e células. 
\indent As evoluções do MFDM tornavam a apresentação da técnica cada vez mais similar à do MEF. Ao atrelar os 
graus de liberdade a funções-teste no elemento, as disparidades em relação à técnica original cresceram a ponto de a formulação se assemelhar 
a uma generalização do FEM. Essas diferenças levaram à documentação, por \citet{basicprinciplesvirtualelementmethods2013a}, de uma nova técnica numérica batizada de Método dos Elementos 
Virtuais (MEV). Para expor a nova metodologia, foi lançada uma série de artigos que apresentavam as principais características do método, bem 
como demonstrações dos critérios de convergência e de unicidade dos graus de liberdade para os exemplos apresentados.\\
\indent No artigo fundacional do método, por motivos de simplificação, tomou-se a equação de Poisson como exemplo de aplicação a equações 
elípticas — sem variação no tempo — e em campos escalares. Subsequentemente, \citet{virtualelementslinearelasticityproblems2013} desenvolveram os princípios do MEV 
aplicado a campos vetoriais, estudando os critérios de convergência e a unicidade dos graus de liberdade voltados a problemas de elasticidade 
linear compressível e quase incompressível.  \\
\indent Embora a aplicação do método em campos vetoriais representasse um avanço significativo, tanto ela quanto a equação de Poisson tratavam 
de elementos com conformidade $H^1$. Isso significa que apenas a variável de campo da equação diferencial deve possuir conformidade no 
elemento. Para expandir esse cenário, \citet{virtualelementmethodsplatebendingproblems2013a} apresentaram os fundamentos e os critérios de convergência do método aplicado ao problema de flexão 
de placas de Kirchhoff-Love, solucionando o problema de conformidade $H^2$, no qual a variação da variável de campo da equação diferencial 
também deve ser conforme no elemento. Buscando generalizar a conformidade do espaço, \citet{virtualelementmethodarbitraryregularity2014} desenvolveram uma família de elementos capaz de se 
adaptar não apenas a malhas não estruturadas, mas também a problemas de regularidade arbitrária $C^{\alpha}$. \\
\indent Na formulação clássica do MEV, a aproximação da solução no elemento por funções polinomiais é realizada com 
a utilização de um projetor na norma $H^p$. Ao trabalhar com a equação em sua forma variacional, o termo-fonte e a matriz de massa (quando 
presente) não apresentam operadores elípticos em suas formas bilineares. A aproximação desses termos por um projetor no espaço $H^p$, apesar 
de aceitável, não era ideal. Esse gargalo impeliu \citet{equivalentprojectorsvirtualelementmethods2013a} a desenvolverem uma expansão da formulação clássica, criando um novo operador de 
projeção $L^2$-ortogonal como um espaço de solução estendido para a aproximação desses termos. A partir da implementação desse novo operador, 
\citet{virtualelementmethodgeneralsecondorderellipticproblemspolygonalmeshes2016} apresentaram uma formulação generalizada para qualquer operador elíptico, e não apenas para o Laplaciano como havia sido feito em \citet{basicprinciplesvirtualelementmethods2013a} e 
\citet{virtualelementmethodgeneralsecondorderellipticproblemspolygonalmeshes2016}. Paralelamente, \citet{hdivhcurlconformingvem2014} ampliou a formulação para espaços $H^1(\text{div})$ e $H^1(\text{rot})\text{-conforme}$, permitindo computar 
equações diferenciais em campos vetoriais com operadores divergente e rotacional, os quais são de grande utilidade em problemas de formulação
mista\footnote{A solução de uma equação diferencial é uma função, como a função dos desloacamento, sendo essa chamada de variável 
de campo. Em algumas formulações, a equação diferencial possui mais de uma variável de campo, por exemplo pressão e temperatura, ou 
deslocamento e tensão de maneira que seja resolvida para duas variáveis de campo simultaneamento. As formulações que apresentam esse
tipo de equação diferencial de duas variáveis de campo são chamadas de formulações mistas}.  \\
\indent Enquanto os trabalhos anteriores apresentavam as formulações iniciais do método, nenhum deles abordava os aspectos de implementação 
prática. Em busca de preencher essa lacuna, \citet{hitchhikersguidevirtualelementmethod2014a} escreveram um guia voltado à equação de Poisson desenvolvida no artigo fundacional do método,
 considerando também as contribuições mais recentes, como a expansão da formulação através da introdução do operador de projeção 
 $L^2$-ortogonal. Este artigo, de grande importância histórica, expõe as características de implementação computacional do método, divergindo 
 dos trabalhos anteriores que focavam estritamente nos aspectos matemáticos. Posteriormente, outros trabalhos foram dedicados à apresentação 
 da implementação computacional. Dentre eles, \citet{virtualelementmethod50linesmatlab2017} expõe considerações relacionadas à programação com um código de exemplo em MATLAB para a 
 equação de Poisson em aproximações lineares, e \citet{chen2020programming} busca aprimorar esse programa por meio da vetorização do código. Por fim, a lacuna das 
 aplicações em MATLAB para formulações de alta ordem e outros operadores foi preenchida em \citet{numericalimplementationhighorder2dvirtualelementmethodmatlab2023a} e \citet{mvemmatlabsoftwarepackagevirtualelementmethods2022}. \\
\indent Formulações não conformes são caracterizadas por possuírem condições "relaxadas" de continuidade no contorno, o que garante maior 
flexibilidade. Aplicações não conformes são de grande utilidade em problemas da mecânica do contínuo, tais como fluxo de fluidos e problemas 
de elasticidade. Essas formulações surgiram inicialmente no MEF, e \citet{nonconformingvirtualelementmethod2016} as adaptou ao MEV 
para elementos bi e tridimensionais, englobando qualquer grau de precisão. Na mesma linha, \citet{conformingnonconformingvirtualelementmethodsellipticproblems2016a} desenvolveu uma estrutura abstrata unificada 
para as formulações conforme e não conforme de maneira que, para esta última, a construção do projetor $L^2$ dependa apenas dos graus de 
liberdade. \\
\indent \citet{illconditioningvirtualelementmethodstabilizationsbases2018} realizou um estudo sobre o mau condicionamento da matriz de rigidez, levando em consideração as bases polinomiais e a 
estabilização utilizada. Em seus resultados, nota-se que as técnicas de estabilização avaliadas apresentam pequenas diferenças entre si 
quanto ao efeito no condicionamento da matriz. Em relação à escolha das bases polinomiais, conclui-se que para problemas de alta ordem ou 
quando os elementos são mal formados — como contornos pequenos ou com formato de estrela não uniforme —, existem outras bases que se adequam 
melhor do que os monômios escalados tradicionalmente utilizados. \citet{exploringhighorderthreedimensionalvirtualelementsbasesstabilizations2018} expandiu esse estudo para problemas em três dimensões. No trabalho, os 
autores propõem variações para as bases nos graus de liberdade das faces/volume e de estabilização, resultando em métodos numericamente mais 
robustos. Com base nos resultados, os autores sugerem diferentes tipos de base a depender do método de refino utilizado, seja ele $h$ ou $p$.\\ 
\indent Assim como o MEF possui a família serendipity de elementos, uma adaptação dessa família para o MEV foi desenvolvida por \citet{serendipitynodalvemspaces2016}. Essa mimetização dos elementos serendipity foi fortemente influenciada pela compatibilidade 
entre os dois métodos. Os elementos virtuais serendipity assemelham-se fortemente aos de elementos finitos, mas mantêm a principal 
característica do MEV: a capacidade de adotar elementos com formatos e número de nós arbitrários, com a exceção de um elemento quadrilateral 
assumindo a forma de um triângulo. \\
\indent Como o MEV não precisa de uma expressão explícita para as funções de base e é diretamente definido no 
espaço físico, \citet{virtualelementmethodcurvededges2019} lançou as bases para o desenvolvimento do MEV com contornos curvos em problemas elípticos. O método apresentado consiste 
na utilização de funções polinomiais no parâmetro $t$ para a aresta curva do elemento. No trabalho, os autores apresentam a construção do 
método, a teoria de convergência, testes numéricos e as técnicas de integração a serem usadas para esses elementos.  \\
\indent Uma peculiaridade da formulação para curvas com funções paramétricas é que o espaço possui funções constantes, mas não lineares. Em 
problemas lineares, por exemplo, por mais que a formulação possua convergência, ela não passa no teste de Irons. Para contornar essa barreira, 
\citet{curvilinearvirtualelements2dsolidmechanicsapplications2020} desenvolveu um espaço de funções aprimorado, com uma construção não convencional que se situa parcialmente em um espaço físico e 
parcialmente mapeada de um espaço de parâmetros. Outra solução, desenvolvida por \citet{polynomialpreservingvirtualelementscurvededges2020}, preserva os polinômios e garante que o resultado não 
dependa da escolha de parametrização da curva. \\
\indent Uma característica marcante do MEV é a necessidade de utilização de um termo de estabilização. Como 
destacado por \citet{lowestorderstabilizationfreevirtualelementmethod2dpoissonequation2025}, o termo de estabilização possui efeitos tanto na estabilidade quanto no condicionamento do método, mas permanece sendo 
um componente escolhido de maneira arbitrária. Buscando superar esse dilema, os autores desenvolveram uma formulação melhorada do método 
utilizando projeções polinomiais de maior ordem; a modificação proposta remove a necessidade do uso do termo de estabilização. Em problemas 
de elasticidade, a extensão dessa formulação livre de estabilização foi realizada por \citet{lowestorderstabilizationfreevirtualelementmethod2dpoissonequation2025}. Visando evitar a utilização de graus de liberdade 
internos em aproximações de maior ordem, \citet{stabilizationfreeserendipityvirtualelementmethodplaneelasticity2023} estenderam o MEV sem estabilização para a família serendipity de 
elementos virtuais. \\
\indent Mesmo que a solução numérica do MEV pertença a um espaço contínuo, ela não pode ser acessada sem a 
resolução de uma equação diferencial parcial em cada elemento polítopo — como será visto em seções posteriores. Desse modo, apenas uma 
projeção de um espaço polinomial descontínuo é acessível, resultando em uma solução descontínua. Outra dificuldade encontrada no método é a 
introdução de um termo de estabilização inerentemente isotrópico; esse termo injeta uma medida de isotropia na solução discreta que, ao lidar 
com problemas anisotrópicos, polui os resultados. Para contornar isso, \citet{reducedbasisstabilizationpostprocessingvirtualelementmethod2024a} propõem o uso do método das bases reduzidas\footnote{O Método de 
Bases Reduzidas (Reduced Basis Method - RBM) é uma técnica para acelerar simulações numéricas complexas. Ele substitui equações pesadas por 
um modelo simplificado "mais leve", construído a partir de poucas soluções conhecidas do problema, permitindo obter respostas rápidas mantendo 
boa precisão.} para avaliar localmente as funções de base de maneira computacionalmente barata. A técnica de bases reduzidas é eficiente na 
solução de equações diferenciais parciais paramétricas em que é necessário solucionar diversas vezes uma mesma equação para diferentes 
parâmetros. Ao considerar a equação diferencial parcial que define as funções de base do MEV como uma equação parametrizada de um elemento 
fixo, toma-se a geometria poligonal do elemento como parâmetro. Essa ferramenta pode ser usada no pós-processamento para a obtenção de uma 
função $H^1$ — considerando um problema de segunda ordem —, útil para visualização e para analisar valores em pontos internos. Outra 
possibilidade da utilização conjunta das bases reduzidas e do MEV é a reconstrução das funções de base para o desenvolvimento de um termo de 
estabilização mais adequado ao problema a ser solucionado.  \\
\indent Em problemas com domínios e/ou dados não suaves, as soluções de problemas elípticos podem conter singularidades fracas. Essas 
singularidades pioram a taxa de convergência de métodos que assumem uma maior regularidade, como o MEF e o 
MEV. Uma das soluções mais amplamente adotadas para o tratamento de descontinuidades e singularidades no domínio é 
a utilização de aproximações enriquecidas baseadas na estrutura da partição da unidade, tais como o MEF Estendidos 
(X-FEM) ou o MEFG. Buscando suplantar as singularidades e descontinuidades, \citet{extendedvirtualelementmethodlaplaceproblemsingularitiesdiscontinuities2019a} generalizaram 
os conceitos dos MEF com partição da unidade para o MEV, criando assim o Método dos 
Elementos Virtuais Estendidos (XMEV). \\
\indent Uma revisão bibliográfica quanto à aproximação de autovalores com o MEV foi realizada por \citet{virtualelementapproximationeigenvalueproblems2022a}. Os autores 
comentam que, para a análise de estabilidade em problemas transientes, é essencial um bom conhecimento das propriedades do esquema de 
discretização espectral. Além disso, a presença de formas de estabilização pode introduzir automodos adicionais artificialmente, tornando 
necessário garantir que eles não poluam a porção do espectro de interesse. Após a apresentação de um esquema prático para problemas de 
autovalores, os autores discutem os efeitos dos parâmetros de estabilização e comentam que certas escolhas de parâmetros fazem com que a 
solução discreta convirja para a contínua com ótima ordem assintótica em refinos $p$ ou $h$. 
\section{O método dos elementos virtuais em problemas de elasticidade.}
\indent As aplicações voltadas a problemas elásticos datam desde o início do desenvolvimento do MEV. No artigo 
publicado por \citet{virtualelementslinearelasticityproblems2013}, são tratados os casos compressíveis e quase incompressíveis em problemas lineares elásticos, apresentando a solução 
numérica pelo MEV em duas dimensões. Assim como na maioria dos artigos iniciais, os autores fazem uma breve 
apresentação do problema e das características do método, focando seus esforços em provar a eficiência e as propriedades de convergência da 
técnica recém-desenvolvida. Uma extensão desse trabalho foi feita por \citet{gain2014virtual} para problemas tridimensionais, apresentando tanto a formulação e 
os aspectos teóricos do método, quanto os seus detalhes de implementação. \\
\indent Como variação ao esquema padrão do MEV, \citet{dualhybridvirtualelementmethodplaneelasticityproblems2020} propõem um esquema duplo híbrido em problemas lineares 
elásticos no estado plano. As variáveis do problema passam a ser tanto os campos de deslocamentos quanto os de tensões. No esquema 
apresentado, as tensões são inicialmente consideradas para satisfazer as equações de equilíbrio localmente em cada elemento, sendo então 
exigido que maximizem a energia complementar local. Outro aspecto é que os campos de deslocamento são utilizados apenas nos contornos entre 
elementos, assumindo a função de um multiplicador de Lagrange para as restrições de continuidade de tração. \\
\indent Costumeiramente, a recuperação dos campos de tensões no MEV é feita utilizando o mesmo operador de projeção 
empregado para a avaliação dos campos de deformações. Em elementos poligonais não triangulares, os deslocamentos podem gerar campos de 
deformações complexos e, por consequência, causar perda de informação nos campos de tensões subjacentes. Buscando suplantar essa barreira, 
\citet{equilibriumbasedstressrecoveryprocedurevem2019} desenvolvem uma adaptação para o MEV da técnica de recuperação por compatibilidade em patches, tradicionalmente 
utilizada junto ao MEF. \\
\indent Indo além da análise determinística, \citet{stochasticvirtualelementmethodsuncertaintypropagationstochasticlinearelasticity2024} apresentam uma variação estocástica do MEV para problemas lineares 
elásticos que possuam propriedades materiais aleatórias e carregamentos estocásticos. No modelo proposto pelos autores, as equações virtuais 
podem ser resolvidas com os mesmos algoritmos utilizados pelo MEF, como a simulação de Monte Carlo, entre outros. 
Além dos algoritmos clássicos, os autores desenvolveram dois métodos numéricos para a solução das equações dos elementos virtuais 
estocásticos: um baseado em uma expansão polinomial caótica e outro fundamentado em uma aproximação intrusiva fraca. 
\subsection{Pequenas deformações:}
\indent O estudo de problemas elásticos não lineares e inelásticos para pequenas deformações foi iniciado em \citet{virtualelementmethodelasticinelasticproblemspolytopemeshes2015a}. O esquema apresentado pelos 
autores permite a inclusão de um algoritmo de leis constitutivas independente da formulação do MEV, similar ao que 
é feito no MEF. Na mesma linha, \citet{arbitraryorder2dvirtualelementspolygonalmeshespartelasticproblem2017a} apresenta uma estrutura matricial para o método voltado à equação da 
elasticidade. Usando a metodologia proposta, o problema pode ser facilmente estendido para casos de não linearidade quando comparado às 
primeiras formulações desenvolvidas. Em contrapartida aos métodos tradicionais, a forma proposta pelos autores realiza a aproximação para a 
deformação ao invés dos deslocamentos.  \\
\indent Dando continuidade ao trabalho desenvolvido por \citet{arbitraryorder2dvirtualelementspolygonalmeshespartelasticproblem2017a}, a extensão para casos inelásticos foi desenvolvida por \citet{arbitraryorder2dvirtualelementspolygonalmeshespartiiinelasticproblem2017a}. Os autores abordam 
três principais casos inelásticos: o primeiro é o de viscosidade isotrópica de Maxwell generalizada; o segundo trata-se da plasticidade 
clássica de von Mises com encruamento (endurecimento) isotrópico linear ou cinemático; e o terceiro aborda o comportamento constitutivo de 
uma liga com memória de forma, modelado por meio de uma abordagem fenomenológica macroscópica. A abordagem da formulação utilizada pelos 
autores permite a utilização de algoritmos constitutivos de "caixa preta", independentes da formulação utilizada para o Método dos Elementos 
Virtuais. \\
\indent Os problemas de não linearidade até então abordados aproximavam a equação diferencial com funções lineares. Embora os resultados 
tivessem sido satisfatórios, \citet{serendipityvirtualelementformulationnonlinearelasticity2019} propuseram um estudo sobre o comportamento do MEV para funções quadráticas em 
casos hiperelásticos compressíveis isotrópicos. Considerando que aproximações de maior ordem no MEV resultam em nós extras e variáveis 
internas no elemento, a abordagem utilizada pelos autores deu-se através da família serendipity de elementos virtuais. Duas estratégias 
diferentes de estabilização foram estudadas pelos autores no escopo do problema abordado. \\
\indent Formulações mistas no regime de pequenas deformações com o uso da formulação variacional de Hellinger-Reissner em duas dimensões foram 
propostas por \citet{stressdisplacementvirtualelementmethodplaneelasticityproblems2017a}. O esquema apresentado é de baixa ordem, com tensores de Cauchy a priori simétricos. O modelo proposto para a solução do 
problema aproxima o campo de deformações com graus de liberdade de tração, enquanto o campo de deslocamentos no polígono consiste, em essência, 
em movimentos de corpo rígido. Enquanto o campo de deslocamentos é aproximado com o uso de funções polinomiais, o campo de deformações 
discreto possui uma construção similar à da velocidade discreta do problema de Stokes. A abordagem em três dimensões foi posteriormente 
elaborada por \citet{threedimensionalhellingerreissnervirtualelementmethodlinearelasticityproblems2020}. \\
\indent Outra proposição de formulação variacional mista foi feita por \citet{huwashizuvariationalapproachselfstabilizedvirtualelements2dlinearelastostatics2023}. No artigo, os autores apresentam o MEVs 
para a formulação de Hu-Washizu. Devido à baixa imposição de compatibilidade, o modelo de deformações é independente do modelo de 
deslocamentos; os autores aproveitam essa liberdade para o desenvolvimento de uma formulação autoestabilizada. Através de testes numéricos, 
demonstra-se que os elementos virtuais propostos são livres de locking, tanto para o caso quase incompressível quanto para elementos altamente 
distorcidos ou com formato não convexo. 
\subsection{Grandes deformações:}
\indent No campo das deformações finitas, dentre os trabalhos iniciais desenvolvidos, cabe citar \citet{basicformulationsvirtualelementmethodvemfinitedeformations2017}. Duas formulações foram propostas para a 
solução do problema: uma formulação mista com dois campos e outra baseada em deslocamento equivalente. A principal característica dessas 
formulações é a utilização da alteração média de volume no elemento, que pode ser computada de maneira exata em elementos virtuais poligonais 
e poliédricos sob qualquer campo de deformação. Outra contribuição foi o desenvolvimento de um novo esquema de estabilização baseado no traço 
do tensor do módulo tangente material. \\
\indent Uma investigação voltada à não linearidade de deformações finitas no regime elastoplástico foi elaborada por \citet{lowordervirtualelementformulationfiniteelastoplasticdeformations2017}, na qual os autores 
desenvolvem uma nova formulação centrada na minimização da expressão da energia incremental. Uma característica marcante do modelo por eles 
desenvolvido é basear-se em uma pseudo-energia de deformação ao invés da forma fraca convencional. A formulação apresentada lida basicamente 
com constantes e vetores, tornando-se matematicamente simples. Outro destaque do trabalho é a introdução de um novo tipo de termo de 
estabilização, envolvendo uma energia de deformação definida positiva modificada.  \\
\indent Na análise de grandes deformações, uma das alternativas de implementação é o uso da abordagem corrotacional. A abordagem 
corrotacional considera os efeitos de grandes deslocamentos e rotações como uma decomposição cinemática dos grandes deslocamentos de corpo 
rígido e das pequenas deformações. \citet{enhancedcorotationalvirtualelementmethodlargedisplacementsplaneelasticity2024} abordam o problema utilizando a formulação sem estabilização inicialmente desenvolvida por \citet{lowestorderstabilizationfreevirtualelementmethod2dpoissonequation2025}, 
explorando uma aproximação livre de divergência. 
\section{O método dos elementos virtuais em problemas de dinâmica}
\indent Em problemas com variação no tempo, um dos primeiros trabalhos a tratar sobre o tema foi o de \citet{virtualelementmethodsparabolicproblemspolygonalmeshes2015a}, no qual apresentaram o método para 
problemas parabólicos tomando como base a equação da difusão dependente do tempo. Os autores desenvolveram uma análise teórica completa para 
o erro entre o problema contínuo e o semidiscreto. As aproximações hiperbólicas foram desenvolvidas por \citet{virtualelementmethodshyperbolicproblemspolygonalmeshes2017} com base nas equações de onda 
dependentes do tempo, empregando uma abordagem idêntica àquela utilizada em problemas parabólicos. Problemas semilineares hiperbólicos foram 
tratados posteriormente por \citet{virtualelementmethodsemilinearhyperbolicproblemspolygonalmeshes2019}. \\
\indent Em elastodinâmica, \citet{antonietti2020virtual} realizaram investigações teóricas e numéricas de desempenho do MEV em problemas 
dependentes do tempo. Os autores avaliam a eficiência do método tanto para malhas de geometrias diversas com refino $h$ quanto para o aumento 
do grau de aproximação polinomial através do refino $p$. Estudos em elementos não convexos foram conduzidos por \citet{nonconvexmesheselastodynamicsusingvirtualelementmethodsexplicittimeintegration2019a}, que aplicaram em seu 
trabalho uma discretização não convexa em domínios bi e tridimensionais. Eles utilizaram um esquema explícito de integração no tempo, cujo 
passo de tempo crítico é aproximado por meio do máximo autovalor local. Destaca-se no trabalho a introdução de um esquema de estabilização 
baseado em uma matriz diagonal. \\
\indent Estudos relacionados à convergência numérica para integrações explícitas no tempo e à estabilidade do MEV 
em elastodinâmica bidimensional e tridimensional foram executados por \citet{numericalrecipeselastodynamicvirtualelementmethodsexplicittimeintegration2020}. Dentre as abordagens propostas pelos autores, destaca-se a 
extensão do MEV para o uso do método B-bar\footnote{O método $\bar{B}$ ($B$-bar) é uma técnica de integração seletiva no Método dos Elementos
Finitos projetada para evitar o bloqueio volumétrico (volumetric locking) em materiais quase incompressíveis. Ele funciona decompondo a matriz 
de deformação-deslocamento $\mathbf{B}$ em partes desviatórias e volumétricas, substituindo esta última por um valor médio reduzido 
($\bar{\mathbf{B}}$) para evitar o enrijecimento artificial do elemento.} voltado ao estudo de sólidos quase incompressíveis. Os autores 
afirmam que é necessária apenas a modificação do termo de estabilização para estabilizar a parte deviatórica da matriz de rigidez. A estimativa 
do passo de tempo crítico é realizada de duas maneiras distintas: o máximo autovalor de um sistema local e o tamanho efetivo do elemento. Dos 
resultados obtidos, destaca-se que o escalonamento diagonal gera massas nodais positivas no MEV — o que não é observado com a técnica de soma 
de linhas (row-sum) em elementos não convexos, mesmo para espaços polinomiais lineares. A aparição de massas negativas junto à técnica de 
soma de linhas foi relacionada a instabilidades numéricas. \\
\indent \citet{antonietti2021arbitrary} propuseram a investigação teórica e numérica do MEV de alta ordem aplicado ao fenômeno de propagação de 
ondas em meios elásticos. Diferente de trabalhos anteriores, os autores realizam a discretização da matriz de rigidez baseando-se na projeção 
ortogonal do gradiente simétrico, e não no projetor elíptico. Os autores avaliam o desempenho do método em relação à convergência, estabilidade e análise de dispersão-dissipação para diversos grupos distintos de malhas. 
\indent Modelagens dinâmicas não lineares para grandes deformações foram abordadas por \citet{virtualelementformulationfinitestrainelastodynamics2021} em elementos bi e tridimensionais não amortecidos, 
utilizando o esquema de Newmark para a integração no tempo. Para o termo de estabilização, a técnica empregada foi a mesma desenvolvida por 
\citet{virtualelementmethodcontact2016}, com a integração sendo realizada em uma submalha triangulada com os mesmos nós da malha original. No trabalho, demonstrou-se que a parte 
dinâmica não necessita de estabilização para garantir a exatidão e a convergência do procedimento, a menos que a análise de autovalores seja 
estritamente necessária. Utilizando a estrutura proposta pelos autores, não é preciso estabilizar a matriz de massa, todavia os autores 
ressalvam que isso só é válido para problemas em que as equações não são dominadas por reações\footnote{Na modelagem de problemas físicos com
equações diferenciais, são combinados três termos: difusivos, advectivos e reativos. Esses termos estão relacionados às derivadas da variável de
campo. Termos reativos são aqueles que não possuem derivada temporal, como o termo de massa em problemas dinâmicos. Assim, problemas dominados
pelo termo reativo são aqueles que a parcela reativa apresenta maior contribuição que difusiva e/ou advectiva}. \\
\indent Em problemas elastoplásticos dinâmicos, \citet{3dmixedvirtualelementformulationdynamicelastoplasticanalysis2021a} adaptou o método para problemas tridimensionais de deformações finitas com o uso de uma 
formulação mista, levando em consideração o caso de incompressibilidade plástica. O comportamento elastoplástico do sólido é modelado com base 
em um esquema de integração numérica (mais precisamente o método de Newmark), com a vantagem de que as matrizes tangentes e de massa são 
combinadas para a solução. Além da formulação pura de deslocamento, a formulação mista de Hu-Washizu também é utilizada. 
\indent Enquanto a maioria dos problemas numéricos dinâmicos realiza a separação do tempo e do espaço, \citet{spacetimeimplementationwaveequationusingvirtualelements2026} utilizaram uma abordagem holística 
de espaço-tempo através do MEV. Os autores destacam que o fato de o MEV possuir flexibilidade na quantidade de nós 
para lidar com o refino da malha torna-se uma grande vantagem. Assim, o método pode ser utilizado de maneira eficiente e com acurácia em 
malhas construídas com continuidade $C^0$ para diferentes densidades de discretização. \\
\indent Em análises lineares dinâmicas explícitas tridimensionais, a presença de tetraedros sliver (elementos muito finos ou 
degenerados) em malhas de elementos finitos diminui dramaticamente o passo crítico de tempo, tornando o progresso da simulação praticamente 
inviável. Para mitigar isso, \citet{virtualelementsagglomeratedfiniteelementsincreasecriticaltimestepelastodynamicsimulations2022} realizaram um estudo no qual utilizaram elementos virtuais na região vizinha ao elemento degenerado de modo
a criar um elemento poliédrico virtual. Como o MEV é compatível com o MEF (Contanto 
que usem o mesmo grau de interpolação), os autores puderam avaliar o efeito da qualidade do formato do poliedro na estimativa do passo crítico 
de tempo em simulações elastodinâmicas explícitas. Problemas de vibração em placas de Kirchhoff-Love foram abordados por \citet{virtualelementmethodvibrationproblemkirchhoffplates2018}. Os autores 
iniciam com uma formulação variacional do problema espectral dependendo apenas do deslocamento transversal da placa. Estabeleceu-se que o 
esquema resultante oferece uma aproximação correta do espectro e provê ótimas estimativas de ordem de erro para as autofunções, além de uma 
ordem dupla para os autovalores, ou seja, o esquema dos autores não cria falsas frequências de vibração atingindo precisão máxima teórica
e com o dobro de velocidade para as freqûencias naturais.\\
\indent Em equações de placas, \citet{adak2022convergence} considerou uma equação de quarta ordem não linear e não local para a demonstração de deformações em 
placas finas e esbeltas. O MEV foi utilizado na discretização da variável espacial, combinando-se a um esquema 
implícito de integração no tempo. Conforme notado pelos autores, a aparição de termos locais diminui a estrutura esparsa do Jacobiano do 
esquema completamente discreto. Para contornar essa perda, os autores introduziram um novo termo e recuperaram a estrutura esparsa do 
Jacobiano. Para diminuir o custo computacional sem comprometer a taxa de convergência, os autores propuseram um esquema totalmente 
linearizado. 
